\section{PRF Details and Miscellaneous}
\label{app:prf}

This appendix records a concrete ZK-friendly PRF candidate compatible with the canonical ARC protocol.  The
construction is Poseidon2-style, but its use as a keyed PRF is an explicit assumption
(Assumption~\ref{ass:poseidon-prf}).

\subsection{PRF specification}
We use a Poseidon2-style permutation with feed-forward and truncation. Let
\(k\in\mathcal K_{\mathsf{key}}\subseteq\Fq^{\ell_{\mathsf{key}}}\) be the secret key encoded by the signed
\(\veck\), and let \(n\in\Fq^{\ell_{\nonce}}\) encode the public nonce. Let
\(t=\ell_{\mathsf{key}}+\ell_{\nonce}\) be the permutation width. Define
\[
  \PRF(k,n) \;=\; \mathrm{Tr}\big(P(k\Vert n) + (k\Vert n)\big),
\]
where \(\mathrm{Tr}:\Fq^{t}\to \Fq^{\ell_{\prftag}}\) outputs the first \(\ell_{\prftag}\) lanes and \(P\) is a
Poseidon2-style permutation~\cite{EPRINT:Poseidon2}.

\subsection{Permutation structure}
Let \(d\ge 3\) be the smallest constant such that \(\gcd(d,q-1)=1\). Fix round counts \(R_F\) (external rounds, even) and
\(R_P\) (internal rounds), round constants, and MDS matrices \(M_E,M_I\). The permutation has the form
\[
  P \;=\; E_{R_F}\circ\cdots\circ E_{R_F/2+1}\circ I_{R_P}\circ\cdots\circ I_{1}\circ E_{R_F/2}\circ\cdots\circ E_1.
\]
Each external round applies the S-box to all lanes, while each internal round applies the S-box to lane \(1\) only.
This structure reduces nonlinear constraints in the post-sign NIZK while retaining diffusion.

\paragraph{External and internal rounds.}
Let \(x=(x_1,\dots,x_t)\in\Fq^t\). For each external round \(i\in[1,R_F]\), fix per-lane constants
\((c^{(i)}_1,\dots,c^{(i)}_t)\in\Fq^t\). For each internal round \(i\in[1,R_P]\), fix a scalar constant \(\hat c^{(i)}\in\Fq\).
Then:
\[
  E_i(x) \;=\; M_E\cdot \bigl((x_1+c^{(i)}_1)^d,\dots,(x_t+c^{(i)}_t)^d\bigr)^\top,
\]
\[
  I_i(x) \;=\; M_I\cdot \bigl((x_1+\hat c^{(i)})^d,\ x_2,\dots,x_t\bigr)^\top.
\]
The matrices \(M_E,M_I\in\Fq^{t\times t}\) are chosen so that they are invertible and provide sufficient diffusion (e.g.,
Cauchy-style MDS matrices as in Poseidon2).

\subsection{Parameter generation}
All PRF parameters are public. A typical parameter generator fixes:
\begin{itemize}[leftmargin=1.25em]
  \item the base field modulus \(q\) (shared with the rest of the protocol),
  \item the S-box exponent \(d\) with \(\gcd(d,q-1)=1\),
  \item width \(t=\ell_{\mathsf{key}}+\ell_{\nonce}\) and truncation \(\ell_{\prftag}\),
  \item round counts \(R_F,R_P\),
  \item MDS matrices \(M_E,M_I\) and round constants \(\{c^{(i)}\},\{\hat c^{(i)}\}\).
\end{itemize}
The truncation length is chosen to satisfy
\(\ell_{\prftag}\ge \lceil \lambda/\log_2(q)\rceil\) as an output-length requirement.  Pseudorandomness is
not implied by this length condition; it is exactly the concrete assumption stated in
Assumption~\ref{ass:poseidon-prf} for \(k\getsunif\mathcal K_{\mathsf{key}}\).

\subsection{Arithmetization strategy (System A with checkpoints)}
\paragraph{Status of the checkpoint schedule.}
Section~\ref{subsec:prf-checkpointed} uses only the abstract interface of three PRF companion families: key binding,
nonlinear checkpoint constraints, and final feed-forward/truncation constraints. The grouped schedule recorded in this
appendix is a candidate realization for the present instantiation.  Any final implementation must fix the exact
round counts, checkpoint positions, and resulting effective degree used in the SmallWood soundness calculation.

To keep showing proofs compact, we commit only selected nonlinear checkpoints:
\begin{itemize}[leftmargin=1.25em]
  \item \textbf{Key lanes.} We commit one row per key lane \(\mathsf{Key}[i]\) and enforce equality to an encoding of
  signed \(\veck\) using selector constraints at fixed \(\Omega\)-coordinates.
  \item \textbf{S-box outputs.} We commit only the S-box outputs at a checkpoint schedule (e.g., every internal round and
  periodically among external rounds). Linear wiring between checkpoints (add-round-constants, MDS multiplications) is
  recomputed by the verifier from opened evaluations and public constants.
\end{itemize}
This yields a System-A style constraint system where the dominant degree comes from composing at most a small number of
unchecked S-box layers between checkpoints.
